% ============================================================================
%  SECCIÓN 4.3: IMPLEMENTACIÓN DE CLASES / COMPONENTES
% ============================================================================

\section{Implementaci\'on de componentes y hooks}

En el ecosistema de React Native, la l\'ogica de negocio se organiza en
componentes de interfaz y hooks personalizados que encapsulan la l\'ogica
reutilizable. A continuaci\'on se describen los componentes y hooks clave
implementados en el proyecto.

\subsection{Componentes reutilizables}

Los componentes reutilizables se encuentran en el directorio
\textit{src/components/} y proporcionan elementos de interfaz comunes a
m\'ultiples pantallas.

\subsubsection{AuthGuard}

El componente \textit{AuthGuard} es un componente de alto orden que envuelve
a los componentes hijos y verifica si el usuario est\'a autenticado. Si no
existe una sesi\'on activa, redirige al usuario a la pantalla de bienvenida:

\begin{lstlisting}
export function AuthGuard({ children }: { children: React.ReactNode }) {
  const { session, loading } = useAuth();
  const router = useRouter();

  useEffect(() => {
    if (!loading && !session) {
      router.replace('/');
    }
  }, [session, loading]);

  if (loading) {
    return <LoadingSpinner message="Verificando sesion..." />;
  }

  return session ? <>{children}</> : null;
}
\end{lstlisting}

\subsubsection{ScenarioCard}

El componente \textit{ScenarioCard} renderiza una tarjeta con la informaci\'on
resumida de un escenario, incluyendo imagen, nombre, tipo y direcci\'on.
Se utiliza tanto en la pantalla de cat\'alogo como en la de favoritos:

\begin{lstlisting}
interface ScenarioCardProps {
  scenario: ScenarioWithDetails;
  onPress: (id: string) => void;
}

export function ScenarioCard({ scenario, onPress }: ScenarioCardProps) {
  const primaryImage = scenario.scenario_images?.find(
    (img) => img.is_primary
  );

  return (
    <TouchableOpacity
      style={styles.card}
      onPress={() => onPress(scenario.id)}
      activeOpacity={0.7}
    >
      {primaryImage && (
        <Image source={{ uri: primaryImage.url }} style={styles.image} />
      )}
      <View style={styles.info}>
        <Text style={styles.name}>{scenario.nombre}</Text>
        <Text style={styles.tipo}>{scenario.tipo}</Text>
        <Text style={styles.direccion} numberOfLines={1}>
          {scenario.direccion}
        </Text>
      </View>
    </TouchableOpacity>
  );
}
\end{lstlisting}

\subsubsection{Componentes de estado}

El proyecto incluye tres componentes para manejar los estados de carga, error
y vac\'io de las listas:

\begin{itemize}
  \item \textit{LoadingSpinner}: muestra un indicador de carga con un mensaje
    descriptivo mientras se obtienen los datos.
  \item \textit{ErrorState}: muestra un mensaje de error con un bot\'on para
    reintentar la operaci\'on fallida.
  \item \textit{EmptyState}: muestra un mensaje informativo cuando no hay
    datos disponibles, con un bot\'on de acci\'on opcional.
\end{itemize}

\subsection{Hooks personalizados}

Los hooks personalizados se encuentran en el directorio \textit{src/hooks/} y
encapsulan la l\'ogica de consulta y mutaci\'on de datos.

\subsubsection{useScenarios}

El hook \textit{useScenarios} gestiona la obtenci\'on de escenarios desde
Supabase. Utiliza \textit{React Query} para el cache y la revalidaci\'on:

\begin{lstlisting}
export function useScenarios() {
  return useQuery({
    queryKey: ['scenarios'],
    queryFn: fetchScenarios,
    staleTime: 1000 * 60 * 5,  // 5 minutos
  });
}

export function useScenario(id: string) {
  return useQuery({
    queryKey: ['scenario', id],
    queryFn: () => fetchScenarioById(id),
    staleTime: 1000 * 60 * 5,
    enabled: !!id,
  });
}
\end{lstlisting}

\subsubsection{useFavorites}

El hook \textit{useFavorites} gestiona las operaciones de consulta,
adici\'on y eliminaci\'on de favoritos. La funci\'on \textit{toggleFavorite}
verifica el estado actual y ejecuta la operaci\'on correspondiente:

\begin{lstlisting}
export function useToggleFavorite() {
  const queryClient = useQueryClient();
  const [isToggling, setIsToggling] = useState(false);

  const toggleFavorite = async (userId: string, scenarioId: string) => {
    setIsToggling(true);
    try {
      const isFav = await checkIsFavorite(userId, scenarioId);

      if (isFav) {
        await removeFavorite(userId, scenarioId);
      } else {
        await addFavorite(userId, scenarioId);
      }

      // Invalidar cache para actualizar la interfaz
      await Promise.all([
        queryClient.invalidateQueries({
          queryKey: ['favorites', userId]
        }),
        queryClient.invalidateQueries({
          queryKey: ['favorite', userId, scenarioId]
        }),
      ]);
    } finally {
      setIsToggling(false);
    }
  };

  return { toggleFavorite, isToggling };
}
\end{lstlisting}
